% !TEX root = ../main.tex
\chapter{Tensor space theory I: over a field}\label{ch:l08}
\begin{paracol}{2}

\begin{sources}
\begin{tabularx}{\linewidth}{@{}l l L@{}}
\toprule
\textbf{Video} & \textbf{Length} & \textbf{Content}\\
\midrule
\seg{8}{1}{Lecture 8: Tensor space theory I: over a field} & 2:22:59 & Review of fields, rings and vector spaces; linear maps, $\mathrm{Hom}$, $\mathrm{End}$, $\mathrm{Aut}$, the dual space; why modules over a non-commutative ring fail; $(p,q)$ tensors, the tensor space, pointwise operations,
the tensor product; $T^0_1V=V^*$ and $T^1_1V\cong\mathrm{End}(V^*)$; Hamel bases and dimension; $V^{**}\cong V$ in finite dimension; basis-free construction; components, dual basis and reconstruction of a tensor; change of basis; the
summation convention and the up/down convention; matrices as ``witchcraft''; transformation of components; the determinant via top forms and volume forms; the determinant of a bilinear form is not well defined; outlook to
principal bundles and tensor densities.\\
\bottomrule
\end{tabularx}

\smallskip
\textbf{Transcript} \file{materials/transcripts/C08-V01.txt} (YouTube automatic captions).
\end{sources}

\begin{reading}[Where this lecture sits in the other notes]
\emph{Gravity and Light}, Lecture~3 (Multilinear algebra); \emph{Theoretische Mechanik}, Lecture~5 (Multilinear algebra). Both cover the tensor space, the dual basis and the change of basis; the extended
treatment of the determinant through top forms and the discussion of matrices as a bad notation are, as far as the titles say, specific to this course; the module generalisation announced here is the subject of \cref{ch:l11}.
(From the lecture titles and the plan of the course; not checked in detail.)
\end{reading}

% ==================================================================
\section{Fields and rings}\label{sec:te-field}

A quick review of vector spaces \ts{8}{1}{0:04}. A vector space is built on a structure called a \emph{field}: an algebraic field, what a mathematician calls a field, not what a physicist calls a field (a function on space-time) \ts{8}{1}{0:25}.

\begin{definition}[New concepts: field]\label{def:te-field}
A \textbf{field} is a triple $(K,+,\cdot)$ where $K$ is a set and
\[
  +\colon K\times K\to K,\qquad \cdot\colon K\times K\to K
\]
are two maps (addition and multiplication) satisfying
\begin{itemize}
  \item for $+$: it is commutative, associative, there is a neutral element (called $0$, whatever element of the set it is), and every element has an inverse with respect to $+$: adding it gives the neutral element
        (the neutral element is to be written down before the inverse, since the inverse refers to it);
  \item for $\cdot$: the same four laws, but only on the set $K\setminus\{0\}$, where the neutral element of addition has been removed: commutativity, associativity, a neutral element (named $1$) and inverses
\end{itemize}
\ts{8}{1}{0:25}\ts{8}{1}{0:48}\ts{8}{1}{1:09}\ts{8}{1}{1:32}\ts{8}{1}{1:53}\ts{8}{1}{2:14}\ts{8}{1}{2:34}.
\end{definition}
(The set of the four laws is written ``CANI'' in the lecture: Commutativity, Associativity, Neutral element, Inverse \ts{8}{1}{1:32}\ts{8}{1}{1:53}.) The multiplicative laws are required only for $K\setminus\{0\}$ because $0$ has no multiplicative inverse. (A student:
could one not just ignore the zero? ``How do you ignore it, you close your eyes when it comes along?'' --- taking it out \emph{is} the meaning of ignoring it \ts{8}{1}{2:34}\ts{8}{1}{2:55}\ts{8}{1}{3:16}.) These are the tools learned in high school \ts{8}{1}{3:16}.
\begin{coursenote}[The distributive law]
The lecture does not mention the law linking $+$ and $\cdot$ ($a\cdot(b+c)=a\cdot b+a\cdot c$), which is also part of the definition of a field (and of a ring); at the board it is left implicit. Not said in class.
\end{coursenote}
\begin{history}[Field and ring]
\emph{Not said in class.} \emph{Checked.} The word \emph{Körper} (field) for a set of numbers closed under the four operations is due to R.~Dedekind (1871); the word \emph{Ring} for the structure without inverses was introduced by D.~Hilbert (\emph{Zahlbericht}, 1897).\par
\emph{References.} R.~Dedekind, supplement to Dirichlet's \emph{Vorlesungen über Zahlentheorie} (1871); D.~Hilbert, ``Die Theorie der algebraischen Zahlkörper'', \emph{Jahresber.\ DMV} 4 (1897).
\end{history}

A weaker notion, ``wildly important'' later when modules are considered, is the following \ts{8}{1}{3:16}\ts{8}{1}{3:39}.
\begin{definition}[New concepts: ring]\label{def:te-ring}
A \textbf{ring} is a set $R$ with two operations $+$ and $\cdot$ of the same shape, but satisfying only: $(R,+)$ is a commutative (abelian) group (CANI for $+$); for $\cdot$ the commutativity is gone, the associativity and a neutral element need to stay (there are rings without neutral
element; one says a \textbf{unital} ring, a ring ``with unit''), and, most importantly, the multiplicative inverses are gone \ts{8}{1}{3:39}\ts{8}{1}{4:01}\ts{8}{1}{4:22}\ts{8}{1}{4:43}.
\end{definition}
Every field is a ring, but there are rings that are not fields \ts{8}{1}{4:43}\ts{8}{1}{5:03}.
\begin{example}[Rings that are not fields]\label{ex:te-rings}
\begin{itemize}
  \item The integers $\mathbb Z$ with $+$ and $\cdot$ (as in first grade): a \emph{commutative} ring, a nice one; but the multiplicative inverse of $2$ is $\tfrac12$, which is not in $\mathbb Z$. One learned to calculate with rings before fields \ts{8}{1}{5:03}\ts{8}{1}{5:24}\ts{8}{1}{5:45}.
  \item $\R$, $\mathbb Q$ and so on, with $+,\cdot$, are fields, ``very special rings'' \ts{8}{1}{5:45}\ts{8}{1}{6:06}.
  \item A \emph{non-commutative} ring: the $M\times M$ matrices over $\R$, which can be added and multiplied by matrix multiplication; the multiplication is not commutative and not every matrix has an inverse (if the determinant is zero)
        \ts{8}{1}{6:06}\ts{8}{1}{6:26}\ts{8}{1}{6:47}. A very natural structure; rings are seen everywhere; function rings later \ts{8}{1}{6:47}\ts{8}{1}{7:08}.
\end{itemize}
\end{example}

% ==================================================================
\section{Vector spaces and linear maps}\label{sec:te-vs}

Now back to fields \ts{8}{1}{7:08}.
\begin{definition}[New concepts: $K$-vector space]\label{def:te-vs}
Let $K$ be a field (important). A \textbf{$K$-vector space} is a triple $(V,\oplus,\odot)$ (the operations are circled because they are \emph{not} the operations of $K$: the $+$ of $K$ is not the $\oplus$ of $V$, and the $\cdot$ of $K$ is not the $\odot$: totally different structure) with
\[
  \oplus\colon V\times V\to V,\qquad \odot\colon K\times V\to V,
\]
the second taking an element of the underlying field and a vector and producing a vector (not two vectors producing a vector). The standard example is a real vector space, over $\R$: vectors can be added and there is a scalar (``$S$'') multiplication. The axioms:
\begin{itemize}
  \item for $\oplus$ the same CANI laws as for a field's addition, now applied to the vector space elements;
  \item for $\odot$ there is no CANI (it does not take two vectors); there are four other laws, linking the operations: an associative law, two distributive laws (one distributes over the $+$ of $K$, the other over the $\oplus$ of $V$), and the unitary law $1\odot v=v$
\end{itemize}
\ts{8}{1}{7:08}\ts{8}{1}{7:29}\ts{8}{1}{7:49}\ts{8}{1}{8:10}\ts{8}{1}{8:30}\ts{8}{1}{8:50}\ts{8}{1}{9:10}\ts{8}{1}{9:31}\ts{8}{1}{37:08}.
\end{definition}
(The axioms are not spelled out at the board: ``if you don't remember this precisely, go back to your linear algebra book and study this now'' \ts{8}{1}{9:31}\ts{8}{1}{9:52}. The list of eight axioms is mentioned again at \ts{8}{1}{37:08}: commutative, associative, neutral, inverse, an associative law, two distributive laws and the unitary law.)

A subset $U\subseteq V$ is a \textbf{vector subspace} if for all $u_1,u_2\in U$ the sum $u_1\oplus u_2$ is again in $U$ and for all $\lambda\in K$ also $\lambda\odot u_1\in U$: $U$ is \emph{closed} under the two operations, they do not lead out of $U$
\ts{8}{1}{9:52}\ts{8}{1}{10:13}\ts{8}{1}{10:34}.

\paragraph{Dropping the circles.} From now on the circles are dropped, $+$ between vectors and $v$ scaled by $\lambda$ as $\lambda v$; one must always know whether it is the plus between elements of the field or between vectors, the product of two elements of $K$ or the scalar multiplication of a $K$ element with a $V$ element:
the expression always tells which, though they are different operations. E.g.\ $\lambda\mu v$ with $\lambda,\mu\in K$ and $v\in V$: the product between $\lambda$ and $\mu$ is the ordinary one in $K$, the others are scalar multiplications \ts{8}{1}{10:34}\ts{8}{1}{10:55}\ts{8}{1}{11:16}\ts{8}{1}{11:38}\ts{8}{1}{12:02}.
The circles are dropped because $\diamond$-type symbols are needed in a moment for another operation \ts{8}{1}{12:02}\ts{8}{1}{12:22}.

\paragraph{Structure preserving maps.} As always one classifies by the structure preserving maps (\cref{int:cl-theme}); for vector spaces the homomorphisms are the \emph{linear maps} \ts{8}{1}{12:22}\ts{8}{1}{12:43}.
\begin{definition}[New concepts: linear map, vector space isomorphism]\label{def:te-linear}
Let $(V,\oplus,\odot)$ and $(W,\boxplus,\boxdot)$ be $K$-vector spaces (different spaces may carry different pluses and scalar multiplications, so they are distinguished here). A map $f\colon V\to W$ is \textbf{linear} if for all $v_1,v_2\in V$ and all $\lambda\in K$
\[
  f(v_1\oplus v_2)=f(v_1)\boxplus f(v_2),\qquad f(\lambda\odot v)=\lambda\boxdot f(v);
\]
one can add first in $V$ and then map, or map the individual vectors and add in the target: the same result (additivity and homogeneity). A bijective linear map is a \textbf{vector space isomorphism}; its inverse is then automatically linear \ts{8}{1}{12:43}\ts{8}{1}{13:04}\ts{8}{1}{13:25}\ts{8}{1}{13:46}\ts{8}{1}{14:07}\ts{8}{1}{14:48}\ts{8}{1}{15:09}.
\end{definition}
(The domain and the target could be the same vector space, operations included, in which case one need not distinguish: but this is the degenerate case, not the rule \ts{8}{1}{14:07}\ts{8}{1}{14:27}.) Two vector spaces $V,W$ (operations suppressed in the notation) are called \textbf{isomorphic as vector spaces} if there is an isomorphism $V\to W$;
``and so on and so on'', it is again an equivalence relation: one considers vector spaces the same if there is a bijective linear map between them, the same idea as always \ts{8}{1}{15:09}\ts{8}{1}{15:30}\ts{8}{1}{15:51}.

% ==================================================================
\section{The vector space of linear maps, End, Aut and the dual}\label{sec:te-hom}

A very important construction \ts{8}{1}{15:51}\ts{8}{1}{16:12}.
\begin{definition}[New concepts: $\mathrm{Hom}(V,W)$]\label{def:te-hom}
$\mathrm{Hom}(V,W)$ is the set of all linear maps from the vector space $V$ to the vector space $W$. In future a tilde on the arrow, $f\colon V\xrightarrow{\ \sim\ }W$, means that the map is linear (``everything will be linear from now on'', it saves chalk). It can again be made into a vector space, by
the operations
\[
  \diamond\!+\ \colon\ (f,g)\mapsto f\mathbin{\diamond\!+}g,\quad (f\mathbin{\diamond\!+}g)(v):=f(v)+g(v),\qquad (\lambda\mathbin{\diamond\!\cdot}f)(v):=\lambda\,f(v),
\]
where the $+$ on the right is the plus of $W$ (the results are added up; first written ``on $V$'', corrected at the board) and the product on the right is the scalar multiplication of $W$: $\mathrm{Hom}(V,W)$ inherits its operations from the target space
\ts{8}{1}{16:12}\ts{8}{1}{16:32}\ts{8}{1}{16:54}\ts{8}{1}{17:16}\ts{8}{1}{17:37}\ts{8}{1}{17:57}\ts{8}{1}{18:19}\ts{8}{1}{18:42}.
\end{definition}
One has to check that $f\mathbin{\diamond\!+}g$ is again a linear map (it is additive; and likewise the definition of $\diamond\cdot$ satisfies the linearity axiom), so that it lies in $\mathrm{Hom}(V,W)$ \ts{8}{1}{18:42}\ts{8}{1}{19:02}\ts{8}{1}{19:23}.

\begin{secondary}[Why a field: modules over a non-commutative ring]
If one had not a field but a ring (removing commutativity and inverses) one could still write down the vector space axioms; this is a wildly different structure which deserves its own name: a \textbf{module over a ring} (to be looked at later in the course, \cref{ch:l11}), nothing but a vector space structure whose underlying
``field'' is a ring. One can do everything above also for modules until one tries to show that $\lambda\diamond\cdot f$, defined by $(\lambda\diamond\cdot f)(v)=\lambda\,f(v)$, is again a linear map: for that one scales, pulls the scaling out, but one cannot pull it in front of the $\lambda$, because
$\lambda$ is a ring element and they no longer commute. So for a module over a non-commutative ring one cannot make $\mathrm{Hom}(V,W)$ again into a module; one can only do it for a module over a \emph{commutative} ring: a very small subtlety that plays a big role
\ts{8}{1}{19:23}\ts{8}{1}{19:44}\ts{8}{1}{20:05}\ts{8}{1}{20:25}\ts{8}{1}{20:46}\ts{8}{1}{21:07}\ts{8}{1}{21:28}.
\end{secondary}
Here everything is fine, and $\mathrm{Hom}(V,W)$ with $\diamond+$ and $\diamond\cdot$ is a vector space in its own right \ts{8}{1}{21:28}\ts{8}{1}{21:49}.

\begin{coursenote}[Vector spaces over different fields]
A student asks whether $V$ and $W$ could be vector spaces over different fields; Schuller first answers: a $K$-vector space $V$ and a $K'$-vector space $W$ can be made into a $K'$-vector space $\mathrm{Hom}(V,W)$, since the scalar multiplication of $\mathrm{Hom}$ is taken
from the target \ts{8}{1}{21:49}\ts{8}{1}{22:10}\ts{8}{1}{22:35}\ts{8}{1}{22:57}\ts{8}{1}{23:18}\ts{8}{1}{23:39}. The same student comes back at \ts{8}{1}{31:42}\ts{8}{1}{32:12}: ``what does linear mean, then?'' and Schuller agrees: the definition of linearity uses the \emph{same} $\lambda$ on both sides, so
\emph{linearity forces $K=K'$}; the first answer fails exactly there (``it was very clever, the question made me consider it'') \ts{8}{1}{32:12}\ts{8}{1}{32:32}\ts{8}{1}{32:53}. So the first answer is wrong as stated and the second one is the correct one.
\end{coursenote}

\paragraph{Terminology.} ``Terrible if you do not know it, trivial once you do'' \ts{8}{1}{23:39}\ts{8}{1}{24:00}.
\begin{definition}[New concepts: endomorphisms, automorphisms, dual space]\label{def:te-end}
\begin{itemize}
  \item $\mathrm{End}(V):=\mathrm{Hom}(V,V)$, the \textbf{endomorphisms} of $V$: the linear maps from $V$ to itself \ts{8}{1}{24:00}\ts{8}{1}{24:21}.
  \item $\mathrm{Aut}(V)$, the \textbf{automorphisms}: the invertible linear maps $V\to V$ (the isomorphisms of $V$ with itself, called isomorphisms before) \ts{8}{1}{24:21}\ts{8}{1}{24:42}.
  \item $V^*:=\mathrm{Hom}(V,K)$, the \textbf{dual vector space} (dual space) of $V$, where $K$ is considered as a $K$-vector space in itself: its addition is the addition of $K$ and its scalar multiplication the multiplication of $K$ \ts{8}{1}{25:44}\ts{8}{1}{26:04}\ts{8}{1}{26:25}.
\end{itemize}
\end{definition}
$\mathrm{Aut}(V)$ is not a vector subspace of $\mathrm{End}(V)$ (the sum of two invertible maps need not be invertible; the lecture says ``vector subspaces'' and then ``usually proper'': there are fewer invertible linear maps than linear maps, and the inclusion $\mathrm{Aut}(V)\subseteq\mathrm{End}(V)$ is proper unless the dimension is zero) \ts{8}{1}{24:42}\ts{8}{1}{25:04}\ts{8}{1}{25:24}\ts{8}{1}{25:44}.
\begin{coursenote}[Aut is not a subspace]
The lecture calls $\mathrm{Aut}(V)$ a vector subspace of $\mathrm{End}(V)$; it is a subset, indeed a group under composition, but not a subspace: it does not contain $0$ and is not closed under addition ($\mathrm{id}+(-\mathrm{id})=0$). Not said in class.
\end{coursenote}
Since $\mathrm{End}(V)$, $\mathrm{Aut}(V)$ and $V^*$ are special cases of $\mathrm{Hom}$, they inherit the operations just constructed (for $V^*$ the operations are $\diamond+$ and $\diamond\cdot$) \ts{8}{1}{26:25}\ts{8}{1}{26:45}\ts{8}{1}{27:06}. The dual space will play an important role: from a vector space and its dual all tensors are constructed \ts{8}{1}{27:06}.

% ==================================================================
\section{Tensors and the tensor space}\label{sec:te-tensors}

\begin{definition}[New concepts: $(p,q)$ tensor]\label{def:te-tensor}
Let $V$ be a $K$-vector space. A \textbf{$(p,q)$ tensor} $T$ over $V$ is a multilinear map
\[
  T\colon \underbrace{V^*\times\dots\times V^*}_{p}\times\underbrace{V\times\dots\times V}_{q}\ \xrightarrow{\ \sim\ }\ K,
\]
a map from $p$ copies of $V^*$ and (Cartesian product) $q$ copies of $V$ to $K$ that is linear in each slot, independently of the other slots: for the linear map one had additivity and scaling, now this applies to each slot
\ts{8}{1}{27:06}\ts{8}{1}{27:28}\ts{8}{1}{27:49}\ts{8}{1}{28:10}.
The pair $(p,q)$ is called the \textbf{valence} of the tensor; sometimes $p+q$ is called its \textbf{rank} \ts{8}{1}{37:29}\ts{8}{1}{37:50}.
\end{definition}
(The name \emph{tensor}: not explained in the lecture; see the coursenote at the end of the section.)

\begin{definition}[New concepts: the tensor space $T^p_qV$]\label{def:te-tspace}
One defines
\[
  T^p_qV:=\underbrace{V\otimes\dots\otimes V}_{p}\otimes\underbrace{V^*\otimes\dots\otimes V^*}_{q}\ :=\ \{\,T\mid T\text{ is a }(p,q)\text{ tensor over }V\,\},
\]
the set of all $(p,q)$ tensors (note that the order of $V$ and $V^*$ in the symbol is the opposite of the arguments in the maps). This is ``just a name'': ``if you ask how this symbol is defined, I don't need to tell you, it is a name; I could also call it $T^p_qV$ identically, an equally good name''
\ts{8}{1}{28:10}\ts{8}{1}{28:32}\ts{8}{1}{28:52}\ts{8}{1}{29:13}\ts{8}{1}{30:38}\ts{8}{1}{30:59}.
It is equipped with the \textbf{pointwise} addition and scalar multiplication: for example for two $(1,1)$ tensors $T,S$ (each eats an element $\sigma\in V^*$ and a vector $v\in V$)
\[
  (T+S)(\sigma,v):=T(\sigma,v)+S(\sigma,v),\qquad(\lambda T)(\sigma,v):=\lambda\,T(\sigma,v),
\]
where the $+$ is the plus in the target $K$, and the product the multiplication in $K$ \ts{8}{1}{29:13}\ts{8}{1}{29:36}\ts{8}{1}{29:57}\ts{8}{1}{30:18}\ts{8}{1}{30:38}.
This space with these operations is the \textbf{$(p,q)$ tensor space over the vector space $V$} \ts{8}{1}{30:59}\ts{8}{1}{31:20}\ts{8}{1}{36:24}\ts{8}{1}{36:44}.
\end{definition}
It is a vector space: one checks the eight axioms (commutative, associative, neutral, inverse, an associative law, two distributive laws and the unitary law) \ts{8}{1}{36:44}\ts{8}{1}{37:08}.

\begin{secondary}[That the tensor space is a set]
The condition that the maps be multilinear is a condition imposed on the maps; to be sure that this is a set one needs the principle of restricted comprehension, so one has to say from which set the elements are taken. Recall that the notation of a map hides that a map is a relation (a functional relation), which is a subset of the Cartesian
product of the domain and the target: so a map $f\colon A\to B$ is an element of the power set $\mathcal P(A\times B)$, and the multilinear maps are selected from the power set of $(V^*)^p\times V^q\times K$. That is how restricted comprehension shows that the $(p,q)$ tensor space is indeed a set (``a general question that I will not
write down every time'') \ts{8}{1}{33:14}\ts{8}{1}{33:37}\ts{8}{1}{33:58}\ts{8}{1}{34:18}\ts{8}{1}{34:39}\ts{8}{1}{35:00}\ts{8}{1}{35:20}\ts{8}{1}{35:41}\ts{8}{1}{36:02}.
\end{secondary}
(See \cref{thm:se-comprehension} for the principle, and \cref{def:cl-map} for maps as relations.)

\begin{definition}[New concepts: tensor product of tensors]\label{def:te-tensorprod}
There are several ways to construct new tensors from old ones; the most important is the \textbf{tensor product}. Unlike the addition, which takes two tensors of the same valence, it takes $T\in T^p_qV$ and $S\in T^r_sV$ of an entirely different valence and produces a tensor of valence $(p+r,q+s)$ over the same $V$:
\[
  (T\otimes S)(\omega_1,\dots,\omega_{p+r},v_1,\dots,v_{q+s}):=T(\omega_1,\dots,\omega_p,v_1,\dots,v_q)\cdot S(\omega_{p+1},\dots,\omega_{p+r},v_{q+1},\dots,v_{q+s}),
\]
where $T$ eats as many co-vectors and vectors as it can, in that order, the product on the right is the multiplication in $K$, and $S$ eats the rest
\ts{8}{1}{37:08}\ts{8}{1}{37:29}\ts{8}{1}{37:50}\ts{8}{1}{38:11}\ts{8}{1}{38:32}\ts{8}{1}{38:52}\ts{8}{1}{39:13}\ts{8}{1}{39:34}\ts{8}{1}{39:55}\ts{8}{1}{40:16}\ts{8}{1}{40:36}.
\end{definition}
\begin{remark}[Two meanings of the symbol $\otimes$]\label{rem:te-two-otimes}
The symbol $\otimes$ is used twice: in $V\otimes\dots\otimes V^*$ as mere notation (the ``orange'' tensor product, a name for the set of \cref{def:te-tspace}), and between two \emph{elements} of tensor spaces, $T\otimes S$, defined as above (the ``blue'' tensor product)
\ts{8}{1}{38:11}\ts{8}{1}{38:32}. The vectors $v_1,\dots,v_p$ here are individual vectors, not components of a tensor or of a vector (components have not been discussed so far) \ts{8}{1}{39:13}\ts{8}{1}{39:34}.
\end{remark}
\begin{coursenote}[Valence indices]
In the lecture the argument of $T$ is called $(V^*)^p\times V^q$ and the tensor space ``$V\otimes\dots\otimes V^*$'' with the order of $V$ and $V^*$ reversed with respect to the arguments; it is the lecture's own remark (\ts{8}{1}{28:32}\ts{8}{1}{28:52}) that this order is different, I keep it so. The name ``tensor''
(Latin \emph{tendere}, to stretch) is not used in class; the derivation is [?] not said in class.
\end{coursenote}

% ==================================================================
\section{First examples: $T^0_1V$ and $T^1_1V$}\label{sec:te-examples}

Examples are in order \ts{8}{1}{40:58}. Start with $T^0_1V$ over a vector space \ts{8}{1}{40:58}\ts{8}{1}{41:19}.
\begin{example}[$(0,1)$ tensors are co-vectors]\label{ex:te-01}
By the notation $T^0_1V=V^*$ (no tensor product appears) and by definition it consists of all multilinear (linear) maps $V\to K$: indeed the definition of $V^*$. The notation is consistent with what was defined before. So the $(0,1)$ tensors are the dual vectors, the whole space is the dual space,
and a $(0,1)$ tensor is called a \textbf{co-vector}, which in physics is its common name \ts{8}{1}{41:19}\ts{8}{1}{41:40}\ts{8}{1}{42:00}\ts{8}{1}{42:21}. What is called a vector and a co-vector depends on what one calls a vector: if the elements of $V$ are called vectors, those of $V^*$ are co-vectors; in finite dimension one can turn this around (below)
\ts{8}{1}{42:21}\ts{8}{1}{42:42}.
\end{example}

\begin{proposition}[$T^1_1V\cong\mathrm{End}(V^*)$]\label{prop:te-11}
The $(1,1)$ tensors over $V$, i.e.\ the multilinear maps $V^*\times V\to K$ (here, as in the lecture, the argument order is $V^*$ first), are isomorphic as a vector space to $\mathrm{End}(V^*)$
\ts{8}{1}{42:42}\ts{8}{1}{43:03}\ts{8}{1}{43:24}\ts{8}{1}{43:44}.
\end{proposition}
\begin{proof}
For every bilinear map construct an element of $\mathrm{End}(V^*)$, and show how to recover precisely the map one started from; then the two contain the same information, one can map one to the other one to one; one must also check that the correspondence is linear, in both directions, which is very simple
\ts{8}{1}{43:44}\ts{8}{1}{44:04}.

Given $T\in T^1_1V$ construct $\hat T\in\mathrm{End}(V^*)$ as follows. $\hat T$ must assign to every $\omega\in V^*$ another element of $V^*$. Use $T$ with its first slot filled with $\omega$: what remains, $v\mapsto T(\omega,v)$, is an object into which one can plug a vector and get a number,
a map $V\to K$, linear since $T$ is linear in the second slot; such a linear map is an element of $V^*$. Define
\[
  \hat T(\omega):=T(\omega,\,\cdot\,)\in V^*
\]
\ts{8}{1}{44:04}\ts{8}{1}{44:26}\ts{8}{1}{44:46}\ts{8}{1}{45:07}\ts{8}{1}{45:28}\ts{8}{1}{45:49}.
No information is lost: $T$ can be reconstructed, since $T(\omega,v)=\bigl(\hat T(\omega)\bigr)(v)$ (the number obtained by letting the co-vector $\hat T(\omega)$ act on $v$), by the very definition of $\hat T$ \ts{8}{1}{45:49}\ts{8}{1}{46:09}\ts{8}{1}{46:29}.
Hence $T\mapsto\hat T$ is a bijection; one checks that it is also linear \ts{8}{1}{46:29}\ts{8}{1}{46:49}. (Linearity of $\hat T(\omega)$ in $\omega$ comes from linearity of $T$ in the first slot, hence $\hat T\in\mathrm{End}(V^*)$; not spelled out in class.)
\end{proof}
One might also feel that $T^1_1V$ is the same as $\mathrm{End}(V)$, but this is generically not true \ts{8}{1}{46:49}\ts{8}{1}{47:10}.

\paragraph{The list of questions.} Continue a list \ts{8}{1}{47:10}\ts{8}{1}{47:30}:
\begin{itemize}
  \item $T^0_1V=V^*$ is the dual.
  \item Any physics textbook says the $(1,0)$ tensors are, as vector space, the same as the vector space $V$ itself, while the $(0,1)$ ones are the duals: is this really true? In general \emph{not} \ts{8}{1}{47:30}\ts{8}{1}{47:51}.
  \item Are the $(1,1)$ tensors the same as $\mathrm{End}(V)$? Yes, but not generically \ts{8}{1}{47:51}\ts{8}{1}{48:11}.
  \item The fundamental question that links the two and answers them: if one takes the dual and then the dual of the dual, does one get $V$ back? Again no, not always, but yes if the vector space is \emph{finite dimensional}
        \ts{8}{1}{48:11}\ts{8}{1}{48:32}.
\end{itemize}
The dimension has not been defined yet: that is what is done next, and only with a notion of dimension can one restrict to the finite dimensional case \ts{8}{1}{48:32}\ts{8}{1}{48:53}.

% ==================================================================
\section{Bases and dimension}\label{sec:te-basis}

Some care is needed in the definition of dimension: here one deals with \emph{bare} vector spaces, no additional structure, no norm. With no further structure than the vector space structure the only notion of basis one can define is the so-called \textbf{Hamel basis}
\ts{8}{1}{48:53}\ts{8}{1}{49:13}\ts{8}{1}{49:35}.
\begin{definition}[New concepts: Hamel basis, dimension]\label{def:te-basis}
Let $(V,+,\cdot)$ be a vector space. A subset $B\subseteq V$ is called a \textbf{Hamel basis} (to avoid confusion with other notions of basis) if
\begin{enumerate}
  \item every finite subset $\{b_1,\dots,b_N\}\subseteq B$ is \emph{linearly independent}, that is
        \[
          \sum_{i=1}^N\lambda^ib_i=0\ \Longrightarrow\ \lambda^1=\dots=\lambda^N=0
        \]
        (the index is put up for future use); and
  \item $B$ is a \emph{generating system}: for every $v\in V$ there exist a finite number $M$ of numbers $v^1,\dots,v^M\in K$ and finitely many elements $b_1,\dots,b_M\in B$ such that $v=\sum_{i=1}^Mv^ib_i$.
\end{enumerate}
The \textbf{dimension} of $V$ is $\dim V:=$ the cardinality of a Hamel basis; it may be finite or infinite (``various infinities'')
\ts{8}{1}{49:35}\ts{8}{1}{49:57}\ts{8}{1}{50:21}\ts{8}{1}{50:45}\ts{8}{1}{51:06}\ts{8}{1}{53:10}\ts{8}{1}{53:31}\ts{8}{1}{53:52}\ts{8}{1}{55:16}\ts{8}{1}{55:37}.
\end{definition}
(That the dimension is well defined, i.e.\ independent of the basis, ``there is something to prove, but we are not going to do that'' \ts{8}{1}{55:16}\ts{8}{1}{55:37}.)
The condition in the first item is a proposition, an implication, either true or false, and its truth or falseness is equivalent to the linear independence of the finite subset; ``this is something that in high school causes students usually great pain to evaluate, mainly because they have not had a proper explanation
of the implication arrow'' (\cref{sec:lo-prop}) \ts{8}{1}{50:45}\ts{8}{1}{51:06}\ts{8}{1}{51:27}\ts{8}{1}{51:47}.

\emph{Why finiteness.} One only looks at \emph{finite} subsets and finite sums. If one required that all the elements of $B$ be linearly independent in the sum sense there would be two problems: $B$ could be uncountable, and then one cannot even write down a sum, not even a series (a series requires partial sums);
and even for a countable basis the sum would become a series up to infinity, but to know what a series does one needs a notion of convergence, a limit; on a vector space without additional structure, not even a topology, one cannot talk about convergence. So on a bare vector space the only thing one can do is to check that for every finite
subset the condition holds \ts{8}{1}{51:47}\ts{8}{1}{52:08}\ts{8}{1}{52:28}\ts{8}{1}{52:49}\ts{8}{1}{53:10}.
\begin{example}[Fourier analysis]\label{ex:te-fourier}
In Fourier analysis the cosines and sines over the real numbers are countably many and form a \emph{countable} basis, but only at the price of having Fourier series, infinite sums: so they do not constitute a Hamel basis; they constitute what is called a \textbf{Schauder basis} [?: ``sha basis'' in the transcript], and for that one needs
more structure, a notion of convergence. A Hamel basis for the functions that are Fourier expandable would be \emph{uncountable}, whereas the Schauder basis is countable; and one uses the Schauder bases in practice
\ts{8}{1}{54:12}\ts{8}{1}{54:35}\ts{8}{1}{54:55}\ts{8}{1}{55:16}.
\end{example}
\begin{history}[Hamel]
\emph{Not said in class.} \emph{Checked.} The notion is named after G.~Hamel, who in 1905 used a basis of $\R$ as a vector space over $\mathbb Q$ to construct discontinuous additive functions; the existence of a Hamel basis of every vector space is equivalent to the axiom of choice (\cref{sec:se-choice}).\par
\emph{References.} G.~Hamel, ``Eine Basis aller Zahlen und die unstetigen L\"osungen der Funktionalgleichung $f(x+y)=f(x)+f(y)$'', \emph{Math.\ Ann.}\ 60 (1905).
\end{history}

\begin{theorem}[Double dual]\label{thm:te-dd}
If $\dim V<\infty$, then $V^{**}:=(V^*)^*$ is isomorphic as a vector space to $V$
\ts{8}{1}{55:57}\ts{8}{1}{56:20}.
\end{theorem}
The proof is asked on the upcoming problem sheet (a vector space isomorphism has to be \emph{constructed}; see the exercise below and \cref{exam:ch08}) \ts{8}{1}{55:57}\ts{8}{1}{56:20}\ts{8}{1}{1:52:57}.
\begin{exercise}[The canonical isomorphism $V\to V^{**}$]\label{ex:te-dd}
Construct a vector space isomorphism $V\to V^{**}$ for finite dimensional $V$, without choosing a basis. (Problem sheet.) The long version is in \cref{exam:ch08}.
\end{exercise}

\paragraph{Vectors and co-vectors are interchangeable in finite dimension.} Hence if one has $V$ and $V^*$, one can identify $V$ with $V^{**}$ for $\dim V<\infty$; one could call the elements of $V$ vectors and those of $V^*$ co-vectors, but $V^*$ is a vector space, so its elements could also be called vectors
and then the elements of $V$ would be the duals of those, the co-vectors. So in finite dimension what is a vector and what is a co-vector can be turned around; but in manifold theory one thinks of tangent spaces as the vector space and over the tangent space one constructs the co-vector space, so there is
the \emph{conventional asymmetry}: one calls these vectors and those co-vectors; but one needs to know that one starts with the space $V$. ``No problem in practice, but once you start thinking about this you might get confused, so I tried to confuse you beforehand; I do not know where
that works pedagogically, but what I said is true'' \ts{8}{1}{56:20}\ts{8}{1}{56:41}\ts{8}{1}{57:01}\ts{8}{1}{57:21}\ts{8}{1}{57:42}.
Once $V^{**}\cong V$ one can redo the construction of \cref{prop:te-11} from $V$ to $V$ and quickly see that $T^1_1V\cong\mathrm{End}(V)$ as well, and likewise the $(1,0)$ tensors are $V^{**}\cong V$: but only for finite dimensional $V$ \ts{8}{1}{57:42}\ts{8}{1}{58:02}\ts{8}{1}{58:25}.
(The lecture comes back to this point in answer to a student's question, \cref{int:te-10}, in the part on transformation of components.)

% ==================================================================
\section{Components, the dual basis and reconstruction}\label{sec:te-components}

So far bases were used only to classify vector spaces into finite and infinite dimensional ones, but bases have another property \ts{8}{1}{58:25}\ts{8}{1}{58:45}. Instead of talking about \emph{the vector}, a rather abstract concept, if one wants to hand over a concrete vector one agrees on a basis and gives its
\textbf{components} $v^1,\dots,v^M$ with respect to that basis; another vector by its components. Bases are used to communicate which precise vector one means, always subject to the choice of basis \ts{8}{1}{58:45}\ts{8}{1}{59:06}\ts{8}{1}{59:26}\ts{8}{1}{59:47}.

\begin{intuition}[Construct without bases]\label{int:te-basisfree}
There is a warning: if one has the components of a vector and starts defining new objects in terms of components, there is the danger of doing \emph{illegal constructions} that are not good in vector space theory, in particular constructions that would take a very different shape if one changed bases. The general rule: always construct
everything \emph{without choice of basis}; only once it is constructed, one may introduce a basis to study which array of numbers is needed to provide the new object. ``Some people say a gentleman only chooses a basis if he really must''; if one has a concrete reason to choose a basis in order to calculate something it is fine,
but constructions should be basis free \ts{8}{1}{59:47}\ts{8}{1}{1:00:07}\ts{8}{1}{1:00:28}\ts{8}{1}{1:00:49}.
\end{intuition}

\begin{definition}[New concepts: components of a tensor]\label{def:te-components}
From now on only \emph{finite dimensional} vector spaces. Let $T$ be a $(p,q)$ tensor over $V$ and $e_1,\dots,e_{\dim V}$ a basis of $V$ (that induces a basis on the tensor space, see below), and $\varepsilon^1,\dots,\varepsilon^{\dim V}$ a basis of $V^*$ (it has the same size as $\dim V^*=\dim V$). The \textbf{components} of $T$ with
respect to these bases are the numbers (in the underlying field, say the real numbers)
\[
  T^{a_1\dots a_p}{}_{b_1\dots b_q}:=T\bigl(\varepsilon^{a_1},\dots,\varepsilon^{a_p},e_{b_1},\dots,e_{b_q}\bigr)\in K.
\]
The tensor is an abstract thing, a multilinear map; its components are numbers, a whole lot of numbers \ts{8}{1}{1:00:49}\ts{8}{1}{1:01:09}\ts{8}{1}{1:01:31}\ts{8}{1}{1:01:53}\ts{8}{1}{1:02:13}\ts{8}{1}{1:02:36}\ts{8}{1}{1:02:57}\ts{8}{1}{1:03:17}.
\end{definition}
(The lecturer writes the basis of $V$ as $B_1,\dots,B_{\dim V}$, then renames it $e$, ``more standard'', and forgets for a moment the second basis \ts{8}{1}{1:02:13}\ts{8}{1}{1:02:36}.)

\begin{definition}[New concepts: dual basis]\label{def:te-dualbasis}
Usually one does not choose independent bases on $V$ and $V^*$ (one could, they are two vector spaces); one chooses on $V^*$ a basis \emph{induced} by the previous choice on $V$, by the condition
\[
  \varepsilon^a(e_b)=\delta^a_b,\qquad\delta^a_b:=\begin{cases}1&a=b\\0&a\ne b\end{cases}
\]
(the $a$-th basis co-vector eats the $b$-th basis vector). This is a linear relation with a unique solution: for every choice of basis on $V$ and this condition one gets a unique basis of $V^*$, called the \textbf{dual basis} (of the dual space); it is ``dual'' because it is a basis of a dual space and
because of this condition. One could choose just any basis of the dual space; ``we choose the dual basis of the dual space'', and from now on bases are always dual bases unless said otherwise \ts{8}{1}{1:03:17}\ts{8}{1}{1:03:37}\ts{8}{1}{1:03:58}\ts{8}{1}{1:04:19}\ts{8}{1}{1:04:40}\ts{8}{1}{1:05:43}\ts{8}{1}{1:06:03}\ts{8}{1}{1:06:24}\ts{8}{1}{1:26:15}.
\end{definition}
Why impose this condition and not another? No particular reason: but ``if you have to make a choice, make sure that you do not make choices all over the place; if one choice can already be used to imply another it is usually most efficient to keep the number of choices minimal''; one could argue that this
choice is made once and for all. It is a very convenient induction principle, since many calculations simplify because there are so many zeros around \ts{8}{1}{1:04:40}\ts{8}{1}{1:05:01}\ts{8}{1}{1:05:21}.

The most important point: once one has a basis, and a tensor with its components, one can \emph{reconstruct the abstract tensor from its components} \ts{8}{1}{1:06:24}\ts{8}{1}{1:06:44}.
\begin{theorem}[Reconstruction of a tensor from its components]\label{thm:te-recon}
Let $V$ be finite dimensional, $e_a$ a basis, $\varepsilon^a$ its dual basis, and $T\in T^p_qV$ with components $T^{a_1\dots a_p}{}_{b_1\dots b_q}$. Then
\[
  T=T^{a_1\dots a_p}{}_{b_1\dots b_q}\cdot e_{a_1}\otimes\dots\otimes e_{a_p}\otimes\varepsilon^{b_1}\otimes\dots\otimes\varepsilon^{b_q},
\]
where (summation convention, see below) the dot is the scalar multiplication in the tensor space $T^p_qV$ (number in the field times a given tensor), the basis vectors $e_a$ are $(1,0)$ tensors and the $\varepsilon^b$ $(0,1)$ tensors, and the $\otimes$ is the tensor product of \cref{def:te-tensorprod} (the ``blue'' one)
\ts{8}{1}{1:06:44}\ts{8}{1}{1:07:07}\ts{8}{1}{1:07:28}\ts{8}{1}{1:07:49}\ts{8}{1}{1:08:10}\ts{8}{1}{1:08:34}\ts{8}{1}{1:08:55}\ts{8}{1}{1:09:16}.
\end{theorem}
Reading: the string $e_{a_1}\otimes\dots\otimes\varepsilon^{b_q}$ is a $(p,q)$ tensor, as is $T$; there are many such tensors, one for each particular choice of $a_1=7,\dots,a_p=28,b_1=1,\dots,b_q=1$ and so on, and one multiplies each by the number $T^{a_1\dots}{}_{b_1\dots}$
\ts{8}{1}{1:08:34}\ts{8}{1}{1:08:55}\ts{8}{1}{1:09:16}.
\begin{proof}
Not carried out in class: one evaluates the right-hand side on an arbitrary collection of co-vectors and vectors, expanding them in the bases, and uses how the $e$'s and the $\varepsilon$'s act on arbitrary vectors and co-vectors; ``it is very basic'' and a finger exercise on the next problem sheet (\cref{exam:ch08});
``this equal sign here is correct because it is in every textbook, and it has practical importance'' \ts{8}{1}{1:09:16}\ts{8}{1}{1:09:38}\ts{8}{1}{1:09:59}\ts{8}{1}{1:10:20}.
\end{proof}
\begin{exercise}[Reconstruction]\label{ex:te-recon}
Verify \cref{thm:te-recon} for a $(1,1)$ tensor by evaluating both sides on $(\varepsilon^c,e_d)$ and on $(\omega,v)=(\omega_a\varepsilon^a,\,v^be_b)$. (Problem sheet; \cref{exam:ch08}.)
\end{exercise}

% ==================================================================
\section{Change of basis and the summation convention}\label{sec:te-change}

Finite dimensional spaces; let $e_1,\dots,e_d$ ($d:=\dim V$) be a basis of whatever vector space \ts{8}{1}{1:10:20}\ts{8}{1}{1:10:41}.
\begin{definition}[Change of basis]\label{def:te-change}
A new basis $\tilde e_1,\dots,\tilde e_d$ (the same number of elements) consists of vectors, each expressible as a linear combination of the old basis (that is what a basis is):
\[
  \tilde e_a=A^b{}_a\,e_b\qquad(\text{sum over }b=1,\dots,d),
\]
with $A^b{}_a\in K$ (two indices: there is a set of components for every new vector $a$), the scalar multiplication being the one of $V$. Not every such combination gives a basis: one could destroy the linear independence and the generating property. The condition is that one can also recover the old basis from the new one,
\[
  e_a=B^b{}_a\,\tilde e_b,
\]
and if both are possible one checks easily that $B$, as a matrix, is the \emph{inverse} of $A$, $B^a{}_mA^m{}_b=\delta^a_b$
\ts{8}{1}{1:10:41}\ts{8}{1}{1:11:02}\ts{8}{1}{1:11:23}\ts{8}{1}{1:11:43}\ts{8}{1}{1:12:05}\ts{8}{1}{1:12:25}\ts{8}{1}{1:12:45}\ts{8}{1}{1:13:07}.
\end{definition}
(Matrices ``in a second''. The index placement of $A$ and $B$ in the transcript is ambiguous; I have chosen the placement that makes the formulas below consistent [?].)

\begin{definition}[New concepts: summation convention]\label{def:te-einstein}
In the reconstruction \cref{thm:te-recon} there are, actually, a succession of $p+q$ sum symbols, $\sum_{a_1=1}^d\dots\sum_{b_q=1}^d$, each running from $1$ to $\dim V$ (the sums over indices of $V$ and of the co-vector space), which were not written; they must be there, since this is the expansion; but by convention they are forgotten. In the future they are not written:
whenever an index appears once up and once down, in the same product (``product in a generalised sense; not if it is in a sum''), it is summed over, unless something else is explicitly said. This is the \textbf{Einstein summation convention}: omit the $\sum_{a=1}^d$ if $a$ appears once up and once down
\ts{8}{1}{1:13:07}\ts{8}{1}{1:13:28}\ts{8}{1}{1:13:49}\ts{8}{1}{1:14:11}\ts{8}{1}{1:14:31}\ts{8}{1}{1:14:52}\ts{8}{1}{1:17:41}\ts{8}{1}{1:18:02}\ts{8}{1}{1:18:23}\ts{8}{1}{1:18:46}.
\end{definition}
\begin{definition}[The up/down convention]\label{def:te-updown}
How does an index deserve to be up or down? There is only one fundamental choice, and everything else follows: the \emph{basis vectors} of $V$ are labelled by a \emph{downstairs} index, $e_1,\dots,e_d$ (each is an element of $V$, not a component: the first vector, the second vector, the $d$-th), and bases of $V^*$ by an \emph{upstairs}
index, $\varepsilon^1,\dots,\varepsilon^d$. Why does everything else follow? In the definition of the components the tensor first eats the co-vectors, which have their index up, so the resulting component has the index up (an equation should match upstairs and downstairs on both sides); the vector slots take $e_b$ with the index down,
so the component has the index down: ``nothing you have to memorise on top; you just have to memorise the fundamental convention''. When one reconstructs, the indices are already in position \ts{8}{1}{1:15:13}\ts{8}{1}{1:15:35}\ts{8}{1}{1:15:57}\ts{8}{1}{1:16:18}\ts{8}{1}{1:16:40}\ts{8}{1}{1:17:01}\ts{8}{1}{1:17:21}\ts{8}{1}{1:17:41}.
\end{definition}
\begin{history}[The summation convention]
\emph{Not said in class.} \emph{Checked.} The convention of omitting the sum sign over a repeated index was introduced by A.~Einstein in his review of general relativity.\par
\emph{References.} A.~Einstein, ``Die Grundlage der allgemeinen Relativit\"atstheorie'', \emph{Ann.\ Phys.}\ 49 (1916).
\end{history}

\begin{remark}[The convention works only with linear maps]\label{rem:te-linear-only}
A remark: one could think ``let us always do this in all branches of mathematics'', a good convention which saves chalk; but it only works with \emph{linear} spaces, in particular linear maps \ts{8}{1}{1:18:46}\ts{8}{1}{1:19:07}\ts{8}{1}{1:19:30}.
Example: let $f$ be a $(1,1)$ tensor, $f\colon V^*\times V\xrightarrow{\sim}K$. Plug in a co-vector $\omega=\omega_a\varepsilon^a$ (basis of $V^*$ with up index, so the components $\omega_a$ have a down index) and a vector $v=v^be_b$:
\[
  f(\omega,v)=f\bigl(\omega_a\varepsilon^a,\,v^be_b\bigr)=\omega_a\,v^b\,f(\varepsilon^a,e_b)=\omega_a\,v^b\,f^a{}_b,
\]
where one used the multilinearity of $f$: additivity pulls out the (invisible) finite sums over $a$ and over $b$, and linearity pulls out the numbers $\omega_a,v^b\in K$ \ts{8}{1}{1:19:30}\ts{8}{1}{1:19:50}\ts{8}{1}{1:20:11}\ts{8}{1}{1:20:31}\ts{8}{1}{1:20:52}\ts{8}{1}{1:21:13}\ts{8}{1}{1:21:34}\ts{8}{1}{1:21:54}\ts{8}{1}{1:22:15}\ts{8}{1}{1:22:35}\ts{8}{1}{1:22:56}.
So the abstract objects $f,\omega,v$ give a number, and one obtains it by building a big sum and product out of their components: one forgets the abstract objects and uses the components with respect to a basis and its dual basis \ts{8}{1}{1:22:56}\ts{8}{1}{1:23:17}\ts{8}{1}{1:23:38}.
\end{remark}
Imagine now one does not see orange (the chalk in which the sums were written in): what is the difference between the first and the second line? Nothing; but the equal sign is not for nothing: it is because of the linearity of $f$ that one can pull out the sum. If one chooses to omit sum signs, an invisible sum can be invisible here or invisible there: ``that is the point
of invisibility, you do not see where it is'', and so the map better be linear, since otherwise it would make a difference where the sum sits, and one would not see it. The summation convention only works when it plays with linear maps; with polynomial maps, as in algebraic geometry, ``and you use the summation convention at least in unclever ways, you are in trouble''
\ts{8}{1}{1:23:38}\ts{8}{1}{1:23:58}\ts{8}{1}{1:24:19}\ts{8}{1}{1:24:39}\ts{8}{1}{1:25:01}\ts{8}{1}{1:25:21}\ts{8}{1}{1:25:42}.

% ==================================================================
\section{Matrices as ``witchcraft mathematics''}\label{sec:te-matrices}

A few polemical remarks about matrices \ts{8}{1}{1:25:42}\ts{8}{1}{1:26:15}. Having chosen a basis in $V$ and the dual basis in $V^*$ it is \emph{tempting} to think of a co-vector $\omega=\omega_a\varepsilon^a$ as the collection of its components and write them as a \emph{row}, $\omega\ \hat=\ (\omega_1,\dots,\omega_d)$, and the vector $v=v^be_b$ as a \emph{column},
$v\ \hat=\ (v^1,\dots,v^d)^{\mathsf T}$; nothing wrong in itself, as the bases are fixed, one has the bases in mind \ts{8}{1}{1:26:15}\ts{8}{1}{1:26:35}\ts{8}{1}{1:26:56}\ts{8}{1}{1:27:17}\ts{8}{1}{1:27:37}\ts{8}{1}{1:27:58}.
Further, $\phi\in\mathrm{End}(V)$ is, as shown, a $(1,1)$ tensor over $V$, $\phi=\phi^a{}_b\,e_a\otimes\varepsilon^b$, and one is tempted, as is the basis of any engineering mathematics, to think of the linear map as a square array with $\phi^1{}_1$ in the top left, the first row containing $\phi^1{}_1,\phi^1{}_2,\dots,\phi^1{}_d$,
then $\phi^2{}_1$, and so on down to $\phi^d{}_d$ in the last line: the \emph{row-column rule}, the first index the row and the second the column \ts{8}{1}{1:27:58}\ts{8}{1}{1:28:18}\ts{8}{1}{1:28:40}\ts{8}{1}{1:29:01}\ts{8}{1}{1:29:21}\ts{8}{1}{1:29:42}\ts{8}{1}{1:30:03}.
(The word ``tempting'' in a mathematics lecture means: be careful \ts{8}{1}{1:26:15}\ts{8}{1}{1:26:35}.)

\begin{intuition}[A matrix is one way to arrange numbers]\label{int:te-matrix}
This is just \emph{one way} to arrange the information, the $d^2$ numbers, as much as a row or column is one way to arrange $d$ numbers: pure convention. A joke on arrangements: if you look for your great-great-grandfather in the cemetery (``I choose this because you have no emotional attachment to him''), he lies in row 7, column 5; if
another scheme were used he would lie elsewhere \ts{8}{1}{1:30:03}\ts{8}{1}{1:30:26}\ts{8}{1}{1:30:46}\ts{8}{1}{1:31:07}. (``Symmetry'' in the transcript is the speech recogniser's rendering of \emph{cemetery}; I read it so throughout.)
\end{intuition}
What does the convention imply? Take two endomorphisms $\phi,\psi\in\mathrm{End}(V)$ and the composition $\phi\circ\psi$, again an endomorphism (the composition of linear maps is linear) \ts{8}{1}{1:31:07}\ts{8}{1}{1:31:28}\ts{8}{1}{1:31:48}.
\begin{proposition}[Matrix multiplication is the representation of composition]\label{prop:te-matmult}
The components of $\phi\circ\psi$ are $(\phi\circ\psi)^a{}_b=\phi^a{}_m\,\psi^m{}_b$, so that in the row-column arrangement they are the matrix product (row of the first times column of the second).
\end{proposition}
\begin{proof}
The component of an endomorphism: $(\phi\circ\psi)^a{}_b=\varepsilon^a\bigl(\phi(\psi(e_b))\bigr)$ (the tensor ``eats'' the result, then $\varepsilon^a$ acts) \ts{8}{1}{1:31:48}\ts{8}{1}{1:32:09}\ts{8}{1}{1:32:30}.
Now $\psi(e_b)=\psi^m{}_b\,e_m$ (quickly checked from what was said before); $\phi$ is linear, so the invisible sum and the numbers $\psi^m{}_b$ are pulled out, leaving $\psi^m{}_b\,\varepsilon^a\bigl(\phi(e_m)\bigr)$, and again $\varepsilon^a$ is linear, so the numbers can be pulled out: $\varepsilon^a(\phi(e_m))=\phi^a{}_m$ is the component of $\phi$. Hence
\[
  (\phi\circ\psi)^a{}_b=\psi^m{}_b\,\phi^a{}_m=\phi^a{}_m\,\psi^m{}_b,
\]
the last step because the product is in $K$, a field, so commutative \ts{8}{1}{1:32:30}\ts{8}{1}{1:32:50}\ts{8}{1}{1:33:10}\ts{8}{1}{1:33:30}\ts{8}{1}{1:33:51}\ts{8}{1}{1:34:12}\ts{8}{1}{1:34:32}.
\end{proof}
Looking through ``cemetery glasses'' one prefers the second order: the indices to be summed over are close together, so that for the first endomorphism the row index and the column index, and for the second again row and column; one sums over the column index of the first and the row index of the second, and gets the standard rule
\emph{row times column}: the entry in row $a$ and column $b$ of $\phi\circ\psi$ is obtained by running over the columns of the first matrix with the row fixed and over the rows of the second with the column fixed \ts{8}{1}{1:34:32}\ts{8}{1}{1:34:53}\ts{8}{1}{1:35:14}\ts{8}{1}{1:35:36}\ts{8}{1}{1:35:57}\ts{8}{1}{1:36:18}\ts{8}{1}{1:36:38}.
Then one only writes the arrays, defines a new rule (the row-times-column rule, \emph{matrix multiplication}) and writes $\hat\phi\,\hat\psi$ for the arrays: ``the whole point of matrix multiplication is that it is the cemetery rule for finding the representing array of the composition. This is the deep reason behind row times column, but it relies on this
additional construction'' \ts{8}{1}{1:36:38}\ts{8}{1}{1:36:59}\ts{8}{1}{1:37:20}.

The rule extends to $\mathrm{Hom}(V,W)$ with corresponding dimensions, and to the action of a co-vector from the left on a vector: without coordinates and bases it is clear that a co-vector acts on a vector to yield a number (this is the definition), and in components
\[
  \omega(v)=\omega_a\varepsilon^a\bigl(v^be_b\bigr)=\omega_a\,v^b\,\varepsilon^a(e_b)=\omega_a\,v^b\,\delta^a_b=\omega_m\,v^m,
\]
where one pulls out the (invisible) sum and the numbers and uses the dual basis; written as a row and a column this is again a matrix product \ts{8}{1}{1:37:20}\ts{8}{1}{1:37:42}\ts{8}{1}{1:38:02}\ts{8}{1}{1:38:23}\ts{8}{1}{1:38:44}.
Without the summation convention $f(\omega,v)=\omega_a\,f^a{}_b\,v^b$: these are all numbers, so they can be written in a different order; one writes $\omega$ as a row, $f$ as a matrix, $v$ as a column, and ends with the expression one has certainly seen: $\omega^{\mathsf T}\,f\,v$, where one writes a \emph{transpose} to emphasise that one comes from the left \ts{8}{1}{1:38:44}\ts{8}{1}{1:39:06}\ts{8}{1}{1:39:27}\ts{8}{1}{1:39:47}.

\begin{remark}[Resist the temptation]\label{rem:te-resist}
The transpose is ``not a good object at all'': very bad notation, almost pretending to be basis independent but not at all. One might think ``$\omega$ comes from the left and $v$ comes from the right, I calculate like in the introductory course''; but that is not necessary, it requires a lot of additional convention and is not worth it. For higher
rank it gets absurd: a $(1,2)$ tensor would have to be a cube of numbers, a $(2,2)$ tensor a four-dimensional cube, ``what is the point?'' It is a singular possibility for $(1,1)$ and $(0,2)$ tensors to come up with such witchcraft notation. One should think in terms of components and the summation rule, and always try to understand
these objects from the component-free, basis-free notation: the task is to gain facility in switching between basis notation and basis-free notation, as done above by plugging in the basis, and to forget the rest, which is confusing. It is mentioned because it is ingrained in our thinking, since it was used often in the beginning; but the lesson
is: do not do this. ``A gentleman only chooses a basis if he must'', and, the stronger principle (the widow principle), ``a lady must never go to the cemetery'': one should enjoy the fortunes of the basis-free notation
\ts{8}{1}{1:39:47}\ts{8}{1}{1:40:08}\ts{8}{1}{1:40:28}\ts{8}{1}{1:40:49}\ts{8}{1}{1:41:09}\ts{8}{1}{1:41:30}\ts{8}{1}{1:41:51}\ts{8}{1}{1:42:12}.
\end{remark}

% ==================================================================
\section{Transformation of components under a change of basis}\label{sec:te-transform}

After this longish aside, the change of basis: define a new basis $\tilde e_a=A^b{}_a\,e_b$ (with the inverse $B$ so that $e_a=B^b{}_a\tilde e_b$) and study how the components of vectors and co-vectors change; the components are basis dependent
\ts{8}{1}{1:42:12}\ts{8}{1}{1:42:32}\ts{8}{1}{1:42:54}.
\begin{proposition}[Transformation of components]\label{prop:te-transform}
Under $\tilde e_a=A^b{}_a\,e_b$, $e_a=B^b{}_a\,\tilde e_b$ (with $B=A^{-1}$), the components of a co-vector $\omega$ and a vector $v$ transform as
\[
  \tilde\omega_a=A^m{}_a\,\omega_m,\qquad \tilde v^a=B^a{}_m\,v^m,
\]
i.e.\ vector components transform with the \emph{inverse} matrix of the one for co-vector components (``vectors transform this way, co-vectors the other way''), and for a $(p,q)$ tensor
\[
  \tilde T^{a_1\dots a_p}{}_{b_1\dots b_q}=B^{a_1}{}_{m_1}\cdots B^{a_p}{}_{m_p}\;A^{n_1}{}_{b_1}\cdots A^{n_q}{}_{b_q}\;T^{m_1\dots m_p}{}_{n_1\dots n_q}:
\]
an upstairs index transforms like a vector index (with $B$), a downstairs index like a co-vector index (with $A$)
\ts{8}{1}{1:42:54}\ts{8}{1}{1:43:15}\ts{8}{1}{1:43:36}\ts{8}{1}{1:43:57}\ts{8}{1}{1:44:18}\ts{8}{1}{1:44:38}\ts{8}{1}{1:44:59}\ts{8}{1}{1:48:26}\ts{8}{1}{1:48:47}\ts{8}{1}{1:49:07}\ts{8}{1}{1:49:28}.
\end{proposition}
\begin{proof}
\emph{Co-vector.} The components are $\tilde\omega_a=\omega(\tilde e_a)$ (applying $\omega$ to the basis vector defines the components: abstract object, plug in the bases, get the components). Using $\tilde e_a=A^m{}_a\,e_m$, linearity (pull out the numbers and the invisible sum), $\tilde\omega_a=A^m{}_a\,\omega(e_m)=A^m{}_a\,\omega_m$
\ts{8}{1}{1:42:54}\ts{8}{1}{1:43:15}\ts{8}{1}{1:43:36}\ts{8}{1}{1:43:57}\ts{8}{1}{1:44:18}.

\emph{Vector.} The $a$-th component of $v$ is the vector applied to the $a$-th co-vector of the dual basis, $v^a=\varepsilon^a(v)$ (the vector considered as a $(1,0)$ tensor eating a co-vector: the identification of $V$ with $V^{**}$ shown above for finite dimension). After a false start (see the coursenote) the lecturer expands $v$ in the
\emph{new} basis, $v=\tilde v^b\tilde e_b$, and expresses the new basis vectors by the old ones, $\tilde e_b=A^m{}_b\,e_m$; then, using the dual basis relation $\varepsilon^a(e_m)=\delta^a_m$ (this is where the dual basis helps) and linearity of $\varepsilon^a$:
\[
  v^a=\varepsilon^a(v)=\tilde v^b\,A^m{}_b\,\varepsilon^a(e_m)=A^a{}_b\,\tilde v^b,
\]
that is $v=A\tilde v$ in components, so $\tilde v^a=B^a{}_m\,v^m$: ``it goes with the opposite matrix, with the inverse of this matrix'' \ts{8}{1}{1:44:59}\ts{8}{1}{1:45:20}\ts{8}{1}{1:45:40}\ts{8}{1}{1:46:01}\ts{8}{1}{1:46:21}\ts{8}{1}{1:46:42}\ts{8}{1}{1:47:03}\ts{8}{1}{1:47:24}\ts{8}{1}{1:47:44}\ts{8}{1}{1:48:05}\ts{8}{1}{1:48:26}.

\emph{Tensors.} Same by multilinearity, slot by slot \ts{8}{1}{1:48:26}\ts{8}{1}{1:48:47}.
\end{proof}
\begin{coursenote}[The first attempt at the vector rule]
At \ts{8}{1}{1:44:59}--\ts{8}{1}{1:47:44} the derivation for the vector components goes wrong at first (an expansion in the wrong basis leads to $\varepsilon^a(\tilde e_m)$, ``I cannot expand that, something is going wrong here''); it is corrected by expanding $v$ in the twiddled basis, as above. The lecturer's names for the matrices
($A$, $B$, the placement of the indices) vary in the transcript and I fixed them as above [?]. The final rule is as stated in class: vectors transform with the inverse matrix with respect to co-vectors.
\end{coursenote}
The matrix $B$ is the inverse of $A$ as a linear map $V\to V$ (``well defined, we know what that is''); $A^a{}_b$ and $B^a{}_b$ are the components of $A$ and $B$ as endomorphisms; ``the inverse matrix'' either uses the matrix multiplication witchcraft or the notion of inverse of a linear map, which is free of coordinates: the relation $A^a{}_mB^m{}_b=\delta^a_b$ \emph{follows} from
$B$ being the inverse, it is not the definition \ts{8}{1}{1:49:28}\ts{8}{1}{1:49:48}\ts{8}{1}{1:50:09}\ts{8}{1}{1:50:29}. This is the transformation rule for a tensor under a change of basis; one has to see where the twiddles, the $A$'s and the $B$'s are, and ``it is best every time you need this, if you want to know explicitly how an upstairs index
transforms, to quickly rederive it, otherwise you get probably confused'' \ts{8}{1}{1:50:29}\ts{8}{1}{1:50:49}\ts{8}{1}{1:51:10}.

A student asks how one goes ``from here to there'' in the step of the vector derivation where $v$ is treated as a $(1,0)$ tensor eating a co-vector \ts{8}{1}{1:51:10}\ts{8}{1}{1:51:31}. The answer:
\begin{intuition}[The $(1,0)$ tensors are not immediately the vectors]\label{int:te-10}
It is immediate that the $(0,1)$ tensors are $V^*$ (all linear maps $V\to K$, by definition). That the $(1,0)$ tensors are $V$ is \emph{not} immediate, and holds only for finite dimension: an element of $T^1_0V$ is a linear map $V^*\to K$, i.e.\ an element of $V^{**}$; to understand a vector as acting on a co-vector one needs
the vector space isomorphism $V\to V^{**}$ of the problem sheet (\cref{thm:te-dd,ex:te-dd}); once one interprets $v\in V$ as an element of $V^{**}$ via this isomorphism, it can act on a co-vector: it is the co-vector acting on the vector, just the other way around. In physics this subtlety is usually not mentioned (``usually this is just presented as: let us say this acts like this'')
\ts{8}{1}{1:51:31}\ts{8}{1}{1:51:55}\ts{8}{1}{1:52:15}\ts{8}{1}{1:52:36}\ts{8}{1}{1:52:57}\ts{8}{1}{1:53:17}\ts{8}{1}{1:53:38}\ts{8}{1}{1:53:59}\ts{8}{1}{1:54:20}\ts{8}{1}{1:54:40}\ts{8}{1}{1:55:01}.
\end{intuition}
The lecturer thanks the student: since the construction of the isomorphism was outsourced to the problem sheet, it could not be understood technically from what was written, ``but it is outsourced there precisely because it is so useful, so you should deal with that'' (a sub-question on the sheet will make it explicit); the construction is not hard, but needs a little more study than the
immediate isomorphism of the $(0,1)$ tensors with the co-vectors \ts{8}{1}{1:53:59}\ts{8}{1}{1:54:20}.

% ==================================================================
\section{Forms, volume and the determinant}\label{sec:te-det}

Something that will be used and that often comes with a lot of confusion: the \textbf{determinant} \ts{8}{1}{1:55:01}\ts{8}{1}{1:55:22}. Trust the order of presentation: first without the basis, then the basis to get the dimension, then back to the basis-free notation, and again introducing components in that precise order; the whole subject must be understood in this order. If one thinks from
the start of a vector as a collection of numbers, its components, one is already on the wrong track: one has to think abstractly and only after introducing a basis is one allowed to do so \ts{8}{1}{1:55:22}\ts{8}{1}{1:55:42}\ts{8}{1}{1:56:02}\ts{8}{1}{1:56:23}.
In ``witchcraft mathematics'' one takes an array of numbers and calculates the determinant according to some mysterious rule, then shows with heavy work that exchanging two rows or two columns gives a minus sign, ``a very mysterious rule acting on a cemetery'' \ts{8}{1}{1:56:23}\ts{8}{1}{1:56:43}.

\paragraph{Two arrays that look alike.} The array also hides a fact: if $f$ is a $(1,1)$ tensor with components $f^a{}_b$ one may arrange them in rows and columns; a $(0,2)$ tensor $g$ has components $g_{ab}$ with both indices down (it eats two vectors, in particular two basis vectors, which have the index down), and these too can be arranged in an array. But an endomorphism
(a $(1,1)$ tensor) and a bilinear form (a $(0,2)$ tensor) are \emph{entirely different beasts}: $f^a{}_b$ has one index up and one down and transforms with an $A$ and a $B$ (the up index with one, the down index with the other) under a change of basis; $g_{ab}$ transforms with two $A$'s. In mathematics one writes such transformations as
$A^{-1}fA$ for the endomorphism and $A^{\mathsf T}gA$ for the bilinear map (``the transpose, as shown before, a concept from witchcraft mathematics''): two totally different objects transform totally differently. If one does not change bases one may delude oneself into thinking that both are matrices and that what can be done with one
can be done with the other, e.g.\ define the determinant ``according to some witchcraft rule'' also for $g_{ab}$: no; \emph{determinants are only defined for endomorphisms}, as seen when done basis free
\ts{8}{1}{1:56:43}\ts{8}{1}{1:57:04}\ts{8}{1}{1:57:26}\ts{8}{1}{1:57:46}\ts{8}{1}{1:58:07}\ts{8}{1}{1:58:28}\ts{8}{1}{1:58:48}\ts{8}{1}{1:59:11}\ts{8}{1}{1:59:31}\ts{8}{1}{1:59:51}\ts{8}{1}{2:00:13}\ts{8}{1}{2:00:33}.

\begin{definition}[New concepts: $n$-form]\label{def:te-form}
Let $V$ be a finite dimensional vector space, $\dim V=d$. An \textbf{$n$-form} is a $(0,n)$ tensor $\omega$ that is \emph{totally anti-symmetric}: for $0<n\le d$,
\[
  \omega(v_1,\dots,v_n)=\operatorname{sgn}(\pi)\;\omega\bigl(v_{\pi(1)},\dots,v_{\pi(n)}\bigr)\qquad\text{for every permutation }\pi\in S_n,
\]
where $S_n$ is the permutation group (``not the sphere'') and $\operatorname{sgn}(\pi)=+1$ for an even and $-1$ for an odd permutation; more simply: if one exchanges any two arguments one picks up a minus sign, exchanging the same or another two picks up yet another minus sign. For $n=0$: a $(0,0)$ tensor is by convention
the underlying number field, a scalar, so a $0$-form is just a number; for $n=1$ a $(0,1)$ tensor the condition has no effect (nothing to exchange), so every co-vector is a $1$-form. $n$ must lie between $0$ and the dimension of the vector space \ts{8}{1}{2:00:33}\ts{8}{1}{2:00:53}\ts{8}{1}{2:01:16}\ts{8}{1}{2:01:38}\ts{8}{1}{2:02:01}\ts{8}{1}{2:02:23}\ts{8}{1}{2:02:44}\ts{8}{1}{2:03:05}\ts{8}{1}{2:03:25}\ts{8}{1}{2:03:46}\ts{8}{1}{2:04:06}.
\end{definition}
(``This anti-symmetry is wildly important.'' For $n>d$ there are anti-symmetric $(0,n)$ tensors in the sense of the definition, but they are zero, since the arguments are then linearly dependent and the anti-symmetry kills everything; so they are not considered \ts{8}{1}{2:04:27}\ts{8}{1}{2:04:47}\ts{8}{1}{2:05:08}.)

\begin{definition}[New concepts: top form, volume form]\label{def:te-top}
A \textbf{top form} on $V$ is a form of the top rank, with as many entries as the vector space has dimensions: a $d$-form \ts{8}{1}{2:04:06}\ts{8}{1}{2:04:27}\ts{8}{1}{2:04:47}\ts{8}{1}{2:05:08}\ts{8}{1}{2:05:28}.
\end{definition}
\begin{proposition}[Top forms are unique up to a factor]\label{prop:te-topuniq}
If $\omega$ and $\omega'$ are two non-vanishing top forms, there is a number $c\in K$ such that $\omega'=c\,\omega$: essentially there is, up to a scalar multiple, only one top form on a vector space \ts{8}{1}{2:05:28}\ts{8}{1}{2:05:49}.
\end{proposition}
\begin{proof}
``Very easy to see because it is totally anti-symmetric: it can only differ by a factor'' \ts{8}{1}{2:05:49}\ts{8}{1}{2:06:10}. Explicitly (not said in class): expanding $d$ vectors $v_i=v_i^{\,a}e_a$ in a basis and using multilinearity and anti-symmetry, a top form is determined by its value on $(e_1,\dots,e_d)$, namely
$\omega(v_1,\dots,v_d)=\bigl(\sum_{\pi\in S_d}\operatorname{sgn}(\pi)\,v_1^{\pi(1)}\cdots v_d^{\pi(d)}\bigr)\,\omega(e_1,\dots,e_d)$, so two top forms are proportional.
\end{proof}
The choice of a top form $\omega$ is a \textbf{choice of volume form} (``choice of volume'') on the vector space: ``the choice is not big, one can choose one, another one is five times the one you chose''. A vector space together with a choice of top form can be called a vector space \emph{with volume}, an additional structure one chooses, all of whose choice lies in the factor
\ts{8}{1}{2:06:10}\ts{8}{1}{2:06:31}\ts{8}{1}{2:06:55}\ts{8}{1}{2:08:39}\ts{8}{1}{2:08:59}.
\begin{definition}[New concepts: volume spanned by vectors]\label{def:te-vol}
Let $v_1,\dots,v_d$ be $d$ vectors in a $d$-dimensional vector space $V$ with a choice of volume form $\omega$. The \textbf{volume spanned} by them (the order matters) is
\[
  \mathrm{vol}(v_1,\dots,v_d):=\omega(v_1,\dots,v_d)
\]
\ts{8}{1}{2:06:55}\ts{8}{1}{2:07:16}\ts{8}{1}{2:07:36}.
\end{definition}
In a three-dimensional space with three vectors, plugging them into the volume form gives a volume; if the three are linearly dependent, if they lie in one plane, the anti-symmetry kills it: the spanned volume is zero; they must be linearly independent to span a non-zero volume, and the same in higher dimensions \ts{8}{1}{2:07:36}\ts{8}{1}{2:07:57}\ts{8}{1}{2:08:18}.
One might have thought that volume needs an inner product, because one needs an inner product to have angles and lengths; but \emph{volume is the choice of a top form}. This will come back later: it is one place where one needs the volume form \ts{8}{1}{2:08:18}\ts{8}{1}{2:08:39}.

\begin{definition}[New concepts: determinant of an endomorphism]\label{def:te-det}
There is no determinant on a vector space without a volume form, ``but it is not that bad'' \ts{8}{1}{2:08:59}\ts{8}{1}{2:09:24}. Let $\phi\in\mathrm{End}(V)$, $V$ finite dimensional of dimension $d$ (emphasised: an endomorphism, a $(1,1)$ tensor, a linear map $V\to V$). Its \textbf{determinant} is the number
\[
  \det\phi:=\frac{\omega\bigl(\phi(e_1),\dots,\phi(e_d)\bigr)}{\omega(e_1,\dots,e_d)}
\]
for some top form $\omega\ne0$ and some basis $e_1,\dots,e_d$ of $V$; ``it is not the determinant of the components, it is a determinant of the abstract object'' \ts{8}{1}{2:09:24}\ts{8}{1}{2:09:44}\ts{8}{1}{2:10:05}\ts{8}{1}{2:10:27}.
\end{definition}
\begin{proposition}[Well-definedness of the determinant]\label{prop:te-detwd}
$\det\phi$ does not depend on the choice of $\omega$ nor on the choice of the basis \ts{8}{1}{2:10:27}\ts{8}{1}{2:10:49}\ts{8}{1}{2:11:12}.
\end{proposition}
\begin{proof}
\emph{Independence of $\omega$} (immediately seen, as the lecturer asks the class): another top form gives a factor $c$ (\cref{prop:te-topuniq}) in the numerator and the same factor $c$ in the denominator, which cancel \ts{8}{1}{2:11:12}\ts{8}{1}{2:11:35}.
\emph{Independence of the basis:} ``you show by plugging in the basis change we had before: you plug it in here, you plug it in there, you pull it out and so on, and you find it is independent of that too, because essentially you get the $\omega$ of the basis change matrix here, and those factors cancel as well; check, not going to do this right here''
\ts{8}{1}{2:11:35}\ts{8}{1}{2:11:55}\ts{8}{1}{2:12:18}.
Explicitly (not said in class): for a new basis $\tilde e_a=A^b{}_a\,e_b$, multilinearity and anti-symmetry give $\omega(\tilde e_1,\dots,\tilde e_d)=\kappa(A)\,\omega(e_1,\dots,e_d)$ with $\kappa(A)=\sum_\pi\operatorname{sgn}(\pi)A^{\pi(1)}{}_1\cdots A^{\pi(d)}{}_d$, and by linearity of $\phi$ likewise
$\omega(\phi(\tilde e_1),\dots,\phi(\tilde e_d))=\kappa(A)\,\omega(\phi(e_1),\dots,\phi(e_d))$; the common factor cancels in the quotient.
\end{proof}
One sees why an endomorphism is needed: the vector $\phi(e_a)$ must again be a vector, to be fed to $\omega$ \ts{8}{1}{2:12:18}.

\paragraph{Reading off the standard rules.} (i) The determinant of $\phi$ coincides with the determinant according to the witchcraft rule applied to the array of its components: ``what you learned is correct, the determinant rule will coincide'' \ts{8}{1}{2:12:18}\ts{8}{1}{2:12:39}. (ii) To what power of the components of $\phi$ does the determinant contribute: each $\phi(e_a)$ gives one
component and $\omega$ eats $d$ of them, hence the products of $d$ components that appear in the witchcraft rule (for a $2\times2$ matrix $ad-bc$, for $3\times3$ ``$a\,d\,m$ minus the other products'' --- the power is the dimension of the vector space) \ts{8}{1}{2:12:39}\ts{8}{1}{2:13:00}\ts{8}{1}{2:13:20}\ts{8}{1}{2:13:41}\ts{8}{1}{2:14:02}.
(iii) Exchanging two rows or two columns gives a minus sign: this can be read off from the definition since $\omega$ is totally anti-symmetric (exchanging in the argument list of $\omega$ the two vectors $\phi(e_a)$ gives a minus sign); all these rules follow directly from the abstract definition \ts{8}{1}{2:14:02}\ts{8}{1}{2:14:23}\ts{8}{1}{2:14:43}\ts{8}{1}{2:15:04}.
(iv) Under a change of basis nothing happens, as already said (independence of the basis): in witchcraft rules $\det(A^{-1}\hat\phi A)=\det\hat\phi$ \ts{8}{1}{2:15:04}\ts{8}{1}{2:15:24}. (v) $\det(\phi\circ\psi)=\det\phi\cdot\det\psi$, the product being in $K$: this is not a matrix product but a composition, the determinant eating an endomorphism, totally basis free; it
follows from the definition by looking at compositions \ts{8}{1}{2:15:24}\ts{8}{1}{2:15:44}\ts{8}{1}{2:16:05}\ts{8}{1}{2:16:26}. ``All the standard rules come up in this way.''

\begin{secondary}[The witchcraft check of independence of the basis]
The lecture also shows the check in matrix language: under a change of basis the representing array changes as $\hat\phi\mapsto A^{-1}\hat\phi A$; the determinant of the product is the product of the determinants, these are numbers, so one can push $\det A^{-1}$ next to $\det A$ and use the rule backwards: $\det A^{-1}\cdot\det A=\det(A^{-1}A)=\det\mathbb 1=1$ (the unit endomorphism, which also gives $1$ from the abstract
definition); so $\det\hat\phi$ does not depend on the basis, ``they better not, we defined them this way''. Not a big result, an important consistency check \ts{8}{1}{2:17:28}\ts{8}{1}{2:17:48}\ts{8}{1}{2:18:08}\ts{8}{1}{2:18:29}\ts{8}{1}{2:18:50}.
\end{secondary}

\paragraph{The determinant of a bilinear form.} In physics, and also in mathematics, one has the opportunity to use the determinant of a bilinear form; this will also be on the problem sheet (\cref{exam:ch08}). As defined the determinant can only be defined for an endomorphism, as one needs to produce vectors again
\ts{8}{1}{2:16:26}\ts{8}{1}{2:16:47}\ts{8}{1}{2:17:07}\ts{8}{1}{2:17:28}. The witchcraft version, to be replaced by a better one on the problem sheet, is to take the determinant of the array $g_{ab}$ of a $(0,2)$ tensor. This cannot coincide with the proper definition, because the object needs to be an endomorphism: it is ``deeply illegal'' from an abstract point of view
\ts{8}{1}{2:18:50}\ts{8}{1}{2:19:10}\ts{8}{1}{2:19:32}.
\begin{proposition}[The determinant of a bilinear form is not well defined]\label{prop:te-detg}
Under a change of basis the components of a $(0,2)$ tensor $g$ transform as $\tilde g=A^{\mathsf T}gA$ (arrays), so that, with the witchcraft determinant,
\[
  \det\tilde g=(\det A)^2\,\det g,
\]
and the determinant of a bilinear map is \emph{not} independent of the choice of basis: it picks up the factor $(\det A)^2$, the square of the determinant of the change of basis matrix
\ts{8}{1}{2:19:32}\ts{8}{1}{2:19:53}\ts{8}{1}{2:20:13}\ts{8}{1}{2:20:34}\ts{8}{1}{2:20:55}.
\end{proposition}
\begin{proof}
$\det(A^{\mathsf T}gA)=\det(A^{\mathsf T})\det g\det A$ and $\det A^{\mathsf T}=\det A$ \ts{8}{1}{2:19:53}\ts{8}{1}{2:20:13}\ts{8}{1}{2:20:34}.
\end{proof}
Hence one would not use it; but it turns out that one does something ill defined from an abstract point of view, and in other places there will be factors coming up in the denominator which are also not well defined under a change of basis, but the \emph{product} of something not well defined and something picking up the opposite factor is independent of the
change of basis: one starts considering these. ``You commit a mistake and later on the opposite mistake and the mistakes cancel: how can we cast this into proper mathematics?'' The answer: by considering \emph{bundles} \ts{8}{1}{2:20:55}\ts{8}{1}{2:21:16}\ts{8}{1}{2:21:37}.
\begin{application}[Principal bundles and tensor densities]\label{app:te-density}
One considers basis changes by taking the base manifold of interest and attaching at every point, as a fibre, the transformation group of the change of bases: this is called a \textbf{principal fibre bundle}; with it one will be able to give a \emph{bundle definition} of what a tensor is, and of what a \textbf{tensor density} is: the density picks up such a factor. All this can be done without hand-waving in the bundle formalism, ``and
that is one of the big reasons why we need bundles'' \ts{8}{1}{2:21:37}\ts{8}{1}{2:21:59}\ts{8}{1}{2:22:19}\ts{8}{1}{2:22:40}. (Principal bundles are treated in \cref{ch:l19}.)
\end{application}

% ==================================================================
\section*{Key concepts}
\begin{keyconcepts}
\textbf{Retrieval questions.}
\begin{enumerate}
  \item State the definition of a field and of a ring. Which axioms does a ring lose, and which example (a) is a ring but not a field, (b) is a non-commutative ring?
  \item Define a $K$-vector space and a linear map. Why can $\mathrm{Hom}(V,W)$ be made into a vector space, and why does the construction fail for modules over a non-commutative ring? Why does linearity force $V$ and $W$ to be over the same field?
  \item Define $\mathrm{End}(V)$, $\mathrm{Aut}(V)$ and $V^*$. Why is $\mathrm{Aut}(V)$ not a subspace?
  \item Define a $(p,q)$ tensor, the tensor space $T^p_qV$ and the tensor product. What are the two meanings of the symbol $\otimes$?
  \item Prove that $T^1_1V\cong\mathrm{End}(V^*)$. Which statements about $T^1_0V$ and $T^1_1V$ hold only in finite dimension?
  \item Define a Hamel basis and the dimension. Why only finite subsets? Why are sines and cosines not a Hamel basis?
  \item Define the components of a tensor, the dual basis, and state the reconstruction theorem.
  \item State the summation convention and the up/down convention. Why does the summation convention need linearity?
  \item How do the components of a vector, a co-vector and a $(p,q)$ tensor change under a change of basis?
  \item Define $n$-forms, top forms and the volume spanned by $d$ vectors; define $\det\phi$ for an endomorphism and show it is well defined. Why is the determinant of a bilinear form not well defined?
\end{enumerate}
\textbf{Concepts.}
\begin{itemize}
  \item Field, ring, vector space, linear map (\cref{def:te-field,def:te-ring,def:te-vs,def:te-linear}); $\mathrm{Hom}$, $\mathrm{End}$, $\mathrm{Aut}$, dual (\cref{def:te-hom,def:te-end}).
  \item $(p,q)$ tensor, tensor space, tensor product (\cref{def:te-tensor,def:te-tspace,def:te-tensorprod,rem:te-two-otimes}); $T^1_1V\cong\mathrm{End}(V^*)$ (\cref{prop:te-11}).
  \item Hamel basis and dimension, double dual (\cref{def:te-basis,thm:te-dd}).
  \item Components, dual basis, reconstruction (\cref{def:te-components,def:te-dualbasis,thm:te-recon}).
  \item Change of basis, summation and up/down conventions, transformation of components (\cref{def:te-change,def:te-einstein,def:te-updown,prop:te-transform}).
  \item Forms, top forms, volume, determinant (\cref{def:te-form,def:te-top,def:te-vol,def:te-det,prop:te-detwd,prop:te-detg}).
\end{itemize}
\end{keyconcepts}

\flushside
\end{paracol}
